\section{Thesis Outline}
\label{sec:intro_outline}

The remainder of the thesis is structured as follows.
Chapter~\ref{chap:theoretical_background} introduces the theoretical concepts needed for continuous, hybrid, and co-simulation models.
Chapter~\ref{chap:requirements} states the user requirements, derives the system requirements, and motivates the choice of Python, \ac{FMI}, OpenSim, and NGSolve as supporting tools.
Chapter~\ref{chap:architecture} presents the \syssimx{} architecture, including the component abstraction, port system, connection model, and master-algorithm interface.
Chapter~\ref{chap:implementation} describes the implementation of structural analysis, algebraic-loop handling, continuous and hybrid master algorithms, and runtime model switching.
It also verifies the main framework mechanisms in focused scenarios.
Chapter~\ref{chap:case_study} evaluates the framework on the controlled-pendulum case study.
The evaluation compares selected co-simulation results with monolithic OpenModelica references and benchmarks the runtime benefit of the multi-model configuration.
Chapter~\ref{chap:discussion_conclusion} discusses the results, states the limitations, draws the conclusions, and outlines future work.